% 第 11 章 Hierarchical Implementation
\bichapter{层次化实现}{Hierarchical Implementation}
\label{chap:hier}

本章说明如何在层次化设计（hierarchical design）上执行综合（synthesis）、布局（placement）、优化（optimization）、时钟树综合（clock tree synthesis）、布线（routing）以及后布线优化（postroute optimization）。内容涵盖：
\begin{itemize}
  \item 使用 Abstract 进行层次化综合
  \item 使用 ETM 进行层次化综合
  \item 层次化综合流程的预检查
  \item 提供块级测试模型（test model）
  \item 指定块级功耗意图（power intent）
  \item Abstract 视图概述
  \item 创建 Abstract 视图
  \item 报告 Abstract 纳入原因
  \item 创建 Abstract 之后修改 block
  \item 创建 Frame 视图
  \item 在顶层链接 Abstract 视图
  \item 将 block 从 Abstract 切换为 Design 视图
  \item 链接带多个 label 的子块
  \item 从顶层指定 block 的可编辑性
  \item 使用 Abstract 做顶层收敛前的准备
  \item 检查含 Abstract 设计的顶层收敛问题
  \item 使用 Abstract 视图执行顶层收敛
  \item 创建 ETM 与 ETM cell library
  \item 在顶层链接 ETM
  \item 使用 ETM 执行顶层收敛
  \item 透明层次优化（Transparent Hierarchy Optimization，THO）
\end{itemize}

% ============================================================
\section{使用 Abstract 进行层次化综合}
\subsection*{Performing Hierarchical Synthesis Using Abstracts}

基于 Abstract 的层次化综合流程主要包括以下步骤：
\begin{enumerate}
  \item 在顶层对设计进行划分（partition），并为下层 block 生成设计库（design library），详见「全芯片设计的划分与规划」。
  \item 对每个下层 block 进行综合，并生成顶层综合所需的块级信息，详见「综合子块」。
  \item 对顶层进行综合，详见「综合顶层」。
\end{enumerate}

\begin{seeAlsoBox}
\begin{itemize}
  \item Overview of Abstract Views（Abstract 视图概述）
  \item Creating Abstract Views（创建 Abstract 视图）
  \item Creating a Frame View（创建 Frame 视图）
\end{itemize}
\end{seeAlsoBox}

\subsection{全芯片设计的划分与规划}
\subsubsection*{Partitioning and Planning the Full Chip Design}

基于 Abstract 的层次化综合流程包含以下两种情形：
\begin{itemize}
  \item 子块与顶层设计的 floorplan 尚不可用；
  \item 子块与顶层设计的 floorplan 已可用。
\end{itemize}

\subsubsection{Floorplan 不可用时的层次化综合流程}
\subsubsection*{Hierarchical Synthesis Flow When Floorplans are not Available}

当子块与顶层设计的 floorplan 尚不可用时，可按以下步骤对顶层做划分与规划，并创建后续自底向上（bottom-up）综合所需的下层 block：

\begin{figure}[htbp]
\centering
\fbox{\parbox{0.85\textwidth}{\centering\vspace{1.5cm}
\textit{（原书 Figure 187：Hierarchical Synthesis Flow When Floorplans are not Available）}\\[0.3em]
\small 请参阅原版 PDF 插图。
\vspace{1.5cm}}}
\caption{Floorplan 不可用时的层次化综合流程}
\label{fig:hier-synth-no-fp}
\end{figure}

\begin{enumerate}
  \item 读入全芯片设计并施加约束，例如：
\begin{lstlisting}
fc_shell> set REF_LIBS "stdcell.ndm macro.ndm"
fc_shell> create_lib TOP -technology techfile \
    -ref_libs $REF_LIBS
fc_shell> analyze -format verilog $rtl_files
fc_shell> elaborate TOP
fc_shell> set_top_module TOP

fc_shell> load_upf fullchip.upf
fc_shell> read_sdc fullchip.sdc
\end{lstlisting}

  \item 识别设计划分并拆分约束，例如：
\begin{lstlisting}
fc_shell> set_budget_options -add_blocks {BLOCK1 BLOCK2}
fc_shell> split_constraints
\end{lstlisting}

  \item 创建带设计信息的子块设计库，并启用按块配置参考库（reference library）的设置，例如：
\begin{lstlisting}
fc_shell> copy_lib -to_lib BLOCK1.nlib -no_design
fc_shell> copy_lib -to_lib BLOCK2.nlib -no_design
fc_shell> set_attribute -objects BLOCK1.nlib \
   -name use_hier_ref_libs -value true
fc_shell> set_attribute -objects BLOCK2.nlib \
   -name use_hier_ref_libs -value true
fc_shell> save_lib -all
\end{lstlisting}

  \item 创建子块设计划分，例如：
\begin{lstlisting}
fc_shell> commit_block -library BLOCK1.nlib BLOCK1
fc_shell> commit_block -library BLOCK2.nlib BLOCK2
fc_shell> save_lib -all
\end{lstlisting}

  \item 加载此前由 \cmd{split\_constraints} 为未映射子块与顶层生成的 UPF 与 SDC 约束，例如：
\begin{lstlisting}
fc_shell> set_constraint_mapping_file ./split/mapfile
fc_shell> load_block_constraints -all_blocks -type SDC  \
   -type UPF -type CLKNET
fc_shell> save_lib -all
\end{lstlisting}

  \item 运行 \cmd{compile\_fusion}，直到子块的逻辑优化（logic optimization）与自动 floorplanning 完成，例如：
\begin{lstlisting}
fc_shell> compile_fusion -to logic_opto
\end{lstlisting}

  \item 运行 \cmd{compile\_fusion}，直到顶层设计的工艺映射（technology mapping）完成，例如：
\begin{lstlisting}
fc_shell> compile_fusion -to initial_map
\end{lstlisting}

  \item 执行设计规划（design planning）操作：从 floorplan 初始化、子块与电压域（voltage area）塑形、硬宏与标准单元布局、电源网络创建，直到 pin assignment。

  \item 运行 \cmd{compile\_fusion}，直到顶层设计的逻辑优化完成，例如：
\begin{lstlisting}
fc_shell> compile_fusion -from logic_opto -to logic_opto
\end{lstlisting}

  \item 仅对顶层做增量标准单元布局，例如：
\begin{lstlisting}
fc_shell> create_placement -floorplan -use_seed_locs
\end{lstlisting}

  \item 执行时序估计（timing estimation）与预算（budgeting），生成块级实现所需的 block budget。

  \item 完成设计规划后，打开 elaborated RTL NDM，加载 floorplan、feedthrough 与 budget，重建子块与顶层设计。
\end{enumerate}

\subsubsection{Floorplan 可用时的层次化综合流程}
\subsubsection*{Hierarchical Synthesis Flow When Floorplans are Available}

当子块与顶层设计的 floorplan 已可用时，可按以下步骤对顶层做划分与规划，并创建后续自底向上综合所需的下层 block：

\begin{figure}[htbp]
\centering
\fbox{\parbox{0.85\textwidth}{\centering\vspace{1.5cm}
\textit{（原书 Figure 188：Hierarchical Synthesis Flow When Floorplans are Available）}\\[0.3em]
\small 请参阅原版 PDF 插图。
\vspace{1.5cm}}}
\caption{Floorplan 可用时的层次化综合流程}
\label{fig:hier-synth-with-fp}
\end{figure}

\begin{enumerate}
  \item 读入全芯片设计并施加约束，例如：
\begin{lstlisting}
fc_shell> set REF_LIBS "stdcell.ndm macro.ndm"
fc_shell> create_lib TOP -technology techfile \
    -ref_libs $REF_LIBS
fc_shell> analyze -format verilog $rtl_files
fc_shell> elaborate TOP
fc_shell> set_top_module TOP

fc_shell> load_upf fullchip.upf
fc_shell> read_sdc fullchip.sdc
\end{lstlisting}

  \item 对设计划分拆分约束，例如：
\begin{lstlisting}
fc_shell> set_budget_options -add_blocks {BLOCK1 BLOCK2}
fc_shell> split_constraints
\end{lstlisting}

  \item 创建带设计信息的子块设计库，并启用按块配置参考库的设置，例如：
\begin{lstlisting}
fc_shell> copy_lib -to_lib BLOCK1.nlib -no_design
fc_shell> copy_lib -to_lib BLOCK2.nlib -no_design
fc_shell> set_attribute -objects BLOCK1.nlib \
   -name use_hier_ref_libs -value true
fc_shell> set_attribute -objects BLOCK2.nlib \
   -name use_hier_ref_libs -value true
fc_shell> save_lib -all
\end{lstlisting}

  \item 创建子块设计划分，例如：
\begin{lstlisting}
fc_shell> commit_block -library BLOCK1.nlib BLOCK1
fc_shell> commit_block -library BLOCK2.nlib BLOCK2
fc_shell> save_lib -all
\end{lstlisting}

  \item 加载此前由 \cmd{split\_constraints} 为未映射子块与顶层生成的 UPF 与 SDC 约束，例如：
\begin{lstlisting}
fc_shell> set_constraint_mapping_file ./split/mapfile
fc_shell> load_block_constraints -all_blocks -type SDC  \
   -type UPF -type CLKNET
fc_shell> save_lib -all
\end{lstlisting}

  \item 加载子块与顶层设计的 floorplan。

  \item 运行 \cmd{compile\_fusion}，直到子块与顶层设计的逻辑优化完成，例如：
\begin{lstlisting}
fc_shell> compile_fusion -to logic_opto
\end{lstlisting}

  \item 执行标准单元布局，随后进行 pin assignment、时序估计与预算。

  \item 完成设计规划后，打开 elaborated RTL NDM，加载 floorplan、feedthrough 与 budget，重建子块与顶层设计。
\end{enumerate}

\begin{noteBox}
各项设计规划操作的细节，请参阅 \textit{Design Planning User Guide}。
\end{noteBox}

\subsection{综合子块}
\subsubsection*{Synthesizing a Subblock}

要综合某个子块并生成顶层综合所需的块级信息，执行以下步骤：
\begin{enumerate}
  \item 读入「全芯片设计的划分与规划」中 rebuild 步骤后生成的子块设计库。
  \item 施加综合该 block 所需的任何块级专用约束与设置。
  \item 使用 \cmd{set\_dft\_signal}、\cmd{set\_scan\_configuration} 等命令施加 DFT 设置，并用 \cmd{create\_test\_protocol} 创建测试协议（test protocol）。
  \item 插入 DFT 逻辑并综合该 block：
\begin{lstlisting}
fc_shell> compile_fusion -to logic_opto
fc_shell> insert_dft
fc_shell> compile_fusion -from initial_place
\end{lstlisting}
  \item 创建只读 Abstract 视图、Frame 视图与测试协议，并保存 block：
\begin{lstlisting}
fc_shell> create_abstract -read_only
fc_shell> create_frame
fc_shell> write_test_model BLOCK1.ctl
fc_shell> save_block -as BLOCK1/PostSynthesis
\end{lstlisting}
\end{enumerate}

\subsection{综合顶层}
\subsubsection*{Synthesizing the Top-Level}

要综合顶层，执行以下步骤：
\begin{enumerate}
  \item 读入「全芯片设计的划分与规划」中 rebuild 步骤后生成的顶层设计。
  \item 设置合适选项并链接顶层设计，例如：
\begin{lstlisting}
fc_shell> set_label_switch_list \
    -reference {BLOCK1 BLOCK2} PostCompile
fc_shell> set_attribute -objects {BLOCK1.nlib BLOCK2.nlib} \
    -name use_hier_ref_libs -value true
fc_shell> set_top_module TOP
\end{lstlisting}
若子块已用 label 保存在设计库中，用 \cmd{set\_label\_switch\_list} 指定要链接的 label；\cmd{set\_top\_module} 用于链接指定的顶层模块。
  \item 施加综合顶层所需的任何专用约束与设置，并读入子块的测试模型，例如：
\begin{lstlisting}
fc_shell> read_test_model BLOCK1.ctl
fc_shell> read_test_model BLOCK2.ctl
fc_shell> create_test_protocol
\end{lstlisting}
  \item 插入 DFT 逻辑并综合顶层设计：
\begin{lstlisting}
fc_shell> compile_fusion -to initial_opto
fc_shell> insert_dft
fc_shell> compile_fusion -from initial_place
\end{lstlisting}
\end{enumerate}

% ============================================================
\section{使用 ETM 进行层次化综合}
\subsection*{Performing Hierarchical Synthesis Using ETMs}

在层次化流程中，block 设计以带 Frame 视图的宏（macro）形式表示；已实现的 block 则以提取时序模型（extracted timing model，ETM）的形式，用于设计的顶层实现。

ETM 是设计的 Liberty 模型表示，由 PrimeTime 工具中的 \cmd{extract\_model} 命令生成。该命令用相关 Liberty 属性捕获设计的时序与功耗信息，示例如下：
\begin{lstlisting}
library("block1") {
...
comment : "PrimeTime Extracted Model." ;
  voltage_map("nom_voltage", 1.08);
  voltage_map("VDD", 1.08);
...
cell( block1 ) {
timing_model_type : "extracted";
...
pg_pin("VDD") {
        voltage_name : "VDD";
        pg_type : primary_power;
}
...
\end{lstlisting}

本节主题包括：
\begin{itemize}
  \item 使用 ETM 运行层次化综合可行性分析流程
  \item 使用 ETM 运行层次化综合流程
  \item 使用 ETM 的层次化综合示例脚本
\end{itemize}

\begin{seeAlsoBox}
\begin{itemize}
  \item SolvNet 文章 2284429，\textit{Top-level Closure Flow With Extracted Timing Models (ETMs)}
  \item Creating ETMs and ETM Cell Libraries（创建 ETM 与 ETM cell library）
  \item Linking to ETMs at the Top Level（在顶层链接 ETM）
\end{itemize}
\end{seeAlsoBox}

\subsection{使用 ETM 运行层次化综合可行性分析流程}
\subsubsection*{Running the Hierarchical Synthesis Feasibility Analysis Flow Using ETMs}

为减少迭代次数并获得更优结果，可在运行层次化综合之前先做可行性分析（feasibility analysis），如图~\ref{fig:hier-etm-feas} 所示。

层次化综合可行性流程使用 \cmd{compile\_fusion} 之后生成的下列块级信息：
\begin{itemize}
  \item Floorplanning 信息，用于步骤 2 中生成 ETM；
  \item 功耗意图与测试模型（CTL）信息，用于步骤 4 的顶层综合。
\end{itemize}

\begin{figure}[htbp]
\centering
\fbox{\parbox{0.85\textwidth}{\centering\vspace{1.5cm}
\textit{（原书 Figure 189：Hierarchical Synthesis Feasibility Analysis）}\\[0.3em]
\small 流程示意：1)~前端工具综合 block（\texttt{compile\_fusion}）；2)~前端工具创建 ETM（\texttt{extract\_model}）；3)~前端工具综合顶层（\texttt{compile\_fusion}）。请参阅原版 PDF 插图。
\vspace{1.5cm}}}
\caption{层次化综合可行性分析}
\label{fig:hier-etm-feas}
\end{figure}

\subsection{使用 ETM 运行层次化综合流程}
\subsubsection*{Running the Hierarchical Synthesis Flow Using ETMs}

按图~\ref{fig:hier-etm-flow} 所述步骤运行层次化综合流程。与可行性流程不同，该流程使用 place and route 之后生成的下列块级信息：
\begin{itemize}
  \item Floorplanning 信息，用于步骤 3 中生成 ETM；
  \item 功耗意图与测试模型（CTL）信息，用于步骤 5 的顶层综合。
\end{itemize}

\begin{figure}[htbp]
\centering
\fbox{\parbox{0.85\textwidth}{\centering\vspace{1.5cm}
\textit{（原书 Figure 190：Hierarchical Synthesis Flow）}\\[0.3em]
\small 流程示意：1)~前端综合 block；2)~后端实现 block；3)~后端创建 ETM；4)~前端综合顶层。请参阅原版 PDF 插图。
\vspace{1.5cm}}}
\caption{层次化综合流程}
\label{fig:hier-etm-flow}
\end{figure}

\subsection{使用 ETM 的层次化综合示例脚本}
\subsubsection*{Example Script of Hierarchical Synthesis Using ETMs}

以下脚本给出层次化综合流程示例：
\begin{lstlisting}
# Specify reference libraries
set REF_LIBS ${std_cell}.ndm
lappend REF_LIBS ${block_design}_etm.ndm
create_lib -tech ${TECH_FILE} -ref_libs $REF_LIBS ${DESIGN_NAME}

# Read the RTL design
analyze -format verilog [list $rtl_list]
elaborate ${DESIGN_NAME}

#ETM.ndm in ref_lib is used to link cell instances
set_top_module ${DESIGN_NAME}

# Load UPF for ETM cell instance
load_upf ${block_design}.upf -scope cell_instance

# Read test model for ETM cell instance
read_test_model ${block_design}.mapped.ctl
...
compile_fusion -to initial_opto

# Query commands
# List all macro cell instances
get_cells -hierarchical -filter "is_hard_macro == true"
# ref view name is "timing" for cell instances bound to ETM
get_attribute ${CELL} is_etm_moded_cell
\end{lstlisting}

% ============================================================
\section{层次化综合流程的预检查}
\subsection*{Performing Prechecks for Hierarchical Synthesis Flows}

要检查层次化综合流程的就绪程度，可使用下列命令：
\begin{itemize}
  \item \cmd{check\_design -checks block\_ready\_for\_top}\\
在 block 综合之后运行，检查已综合的 block 是否已准备好用于顶层综合。
  \item \cmd{check\_design -checks hier\_pre\_compile}\\
在顶层综合之前运行，检查子块与顶层设计是否已准备好进入层次化流程。
\end{itemize}

% ============================================================
\section{提供块级测试模型}
\subsection*{Providing Block-Level Test Models}

在运行层次化综合时，可为顶层设计实现提供子块的测试模型。做法是：在块级综合期间为已映射的 block 生成测试模型，再在顶层综合时将该测试模型注解到子块。例如：

\begin{itemize}
  \item 在块级综合期间生成测试模型：
\begin{lstlisting}
# Read block RTL and constraints
# Specify DFT configuration
set_dft_signal ...
set_scan_configuration ...
create_test_protocol
...
compile_fusion -to initial_opto
insert_dft
write_test_model -output block_des.mapped.ctl
\end{lstlisting}

  \item 在顶层综合中使用块级测试模型：
\begin{lstlisting}
# Read RTL and constraints
# Specify DFT configuration

read_test_model block_des.mapped.ctl
set_dft_signal ...
set_scan_configuration ...
create_test_protocol
...
compile_fusion -to initial_opto
insert_dft
\end{lstlisting}
\end{itemize}

% ============================================================
\section{指定块级功耗意图}
\subsection*{Specifying Block-Level Power Intent}

在运行层次化综合时，必须为顶层设计实现提供子块的功耗意图。可提供完整的块级 UPF，由工具自动提取功耗意图接口。下图示意名为 \texttt{U1} 的子块；顶层功耗意图脚本用 \cmd{load\_upf} 加载块级功耗意图。

\begin{noteBox}
块级功耗意图规范仅在使用 ETM 的层次化综合中必需。在 Abstract 流程中，UPF 约束由工具自动处理。
\end{noteBox}

\begin{itemize}
  \item 在 \texttt{PD\_Top} 设计中将 \texttt{U1} 作为 ETM
  \item 顶层功耗意图：
\begin{lstlisting}
create_power_domain PD_TOP -include_scope
create_supply_net VDD_AO
...
load_upf block.upf -scope U1
...
connect_supply_net VDD_AO -port U1/VDD_AO
\end{lstlisting}
  \item 块级功耗意图：\texttt{block.upf}
\begin{lstlisting}
create_power_domain PD_BLK -include_scope
create_power_switch sw1 -domain PD_BLK \
   -input_supply_port {in VDD_AO} \
   -output_supply_port {out VDD_SW} ...
...
add_port_state sw1/out -state {ON 1.08} -state {OFF off}
...
create_pst pst -supplies {VDD_SW ...}
\end{lstlisting}
\end{itemize}

% ============================================================
\section{Abstract 视图概述}
\subsection*{Overview of Abstract Views}

在 Abstract 视图中，block 的门级网表由仅包含所需接口逻辑（interface logic）的部分门级网表建模；其余逻辑全部移除。

图~\ref{fig:hier-abstract-view} 显示一个 block 及其 Abstract 视图，其中保留了：
\begin{itemize}
  \item 各时序路径上输入端口到第一级寄存器之间的逻辑；
  \item 各时序路径上最后一级寄存器到输出端口之间的逻辑。
\end{itemize}

与纯组合逻辑输入端口到输出端口时序路径（A 到 X）相关的逻辑也会保留。被保留寄存器的时钟连接同样保留。寄存器到寄存器（register-to-register）逻辑则丢弃。

\begin{figure}[htbp]
\centering
\fbox{\parbox{0.85\textwidth}{\centering\vspace{1.5cm}
\textit{（原书 Figure 191：A Block and Its Abstract View）}\\[0.3em]
\small 左侧为完整 Block，右侧为 Abstract View；请参阅原版 PDF 插图。
\vspace{1.5cm}}}
\caption{Block 及其 Abstract 视图}
\label{fig:hier-abstract-view}
\end{figure}

Abstract 视图的接口逻辑包括：
\begin{itemize}
  \item 从输入端口到寄存器或输出端口的时序路径上的所有 cell、pin 与 net；
  \item 从寄存器或输入端口到输出端口的时序路径上的所有 cell、pin 与 net；
  \item 从主时钟（master clock）到生成时钟（generated clock）连接中的任何逻辑；
  \item 驱动接口寄存器的时钟树，包括时钟树中的任何逻辑；
  \item 从时钟端口出发的最长与最短时钟路径；
  \item 属于接口逻辑、且带有时序约束的所有 pin。
\end{itemize}

除接口逻辑外，Abstract 视图还包含与接口逻辑相关的下列信息：
\begin{itemize}
  \item 布局信息（placement information）
  \item 时序约束（timing constraints）
  \item 时钟树例外（clock tree exceptions）
  \item 寄生参数信息（parasitic information）
  \item 悬空 pin（dangling pin）的 transition 与 case value
  \item 功耗管理单元（power management cell）与 UPF 约束
  \item Abstract 视图中保留的 net 上的非默认布线规则（nondefault routing rule）关联，以及最小/最大层约束（若有）
\end{itemize}

% ============================================================
\section{创建 Abstract 视图}
\subsection*{Creating Abstract Views}

在 Fusion Compiler 中，可在设计流程的多个阶段为 block 创建 Abstract 视图，例如综合、布局、优化、时钟树综合、布线之后等。创建 Abstract 视图之前，请确保顶层所需的 scenario 已创建且处于活动（active）状态。

要为顶层收敛创建 Abstract 视图，使用 \cmd{create\_abstract} 命令。
\begin{itemize}
  \item 用 \opt{-timing\_level} 控制 Abstract 视图中时序信息的多少，可选设置如下：
  \begin{itemize}
    \item \texttt{boundary}：创建仅包含连接到各边界端口的边界单元（一层逻辑）的 Abstract。此类 Abstract 还包含 feedthrough 数据路径、feedthrough 组合时钟路径，以及驱动输出端口的内部时钟逻辑。可用于在顶层修复 DRC 违规。
    \item \texttt{compact}：创建紧凑 Abstract，仅包含接口逻辑中关键 setup/hold 时序路径的时序信息。这是默认设置，也是顶层设计收敛的首选。
    \item \texttt{full\_interface}：为全部接口逻辑创建带时序信息的 Abstract。当多重例化块（multiply-instantiated block，MIB）在顶层的约束对各实例不同时，可能需要 \texttt{full\_interface} Abstract。
  \end{itemize}

  \item 要为当前 block 的下层 block 创建 Abstract，可用下列方法之一：
  \begin{itemize}
    \item 用 \opt{-blocks} 并指定 block 名，为特定下层 block 创建 Abstract；
    \item 用 \opt{-all\_blocks} 为所有下层 block 创建 Abstract。
  \end{itemize}
使用 \opt{-blocks} 或 \opt{-all\_blocks} 时，若指定 block 已有 Abstract 视图，默认不会重新创建；要强制重建，使用 \opt{-force\_recreate}。

  \item 要用 \opt{-include\_objects} 把特定 net、port、层次 pin 或 leaf pin 纳入布局或时序 Abstract。

  \item 用 \opt{-read\_only} 创建只读 Abstract。\\
默认情况下，工具创建可编辑（editable）Abstract，可在父 block 上下文中修改；但会把该 block 的 design 视图设为只读，防止对 block 做改动。\\
若创建只读 Abstract，则该 block 的 design 视图仍可编辑，可继续对其修改。

  \item 用 \opt{-preserve\_block\_instances false} 展平物理层次，仅保留下层 block 中相关逻辑。

  \item 用 \opt{-target\_use planning} 或 \opt{-target\_use implementation} 指定 Abstract 用于设计规划还是顶层实现。工具会按预期用途对 Abstract 施加相应内部设置。

  \item 用 \opt{-host\_options} 指定分布式处理的主机选项。当用 \opt{-blocks} 或 \opt{-all\_blocks} 创建多个 Abstract 视图时，可借此缩短运行时间。
\end{itemize}

\subsection{创建含功耗信息的 Abstract}
\subsubsection*{Creating Abstracts With Power Information}

要在块级创建含对应 block 功耗信息的 Abstract，按以下步骤操作：
\begin{enumerate}
  \item 用 \cmd{set\_scenario\_status} 的 \opt{-active true} 激活要进行功耗分析的 scenario，并用下列选项启用功耗分析：
  \begin{itemize}
    \item \opt{-dynamic\_power true}：动态功耗分析
    \item \opt{-leakage\_power true}：漏电功耗分析
  \end{itemize}
功耗信息仅对已启用动态或漏电功耗的活动 scenario 存储。
  \item （可选）通过 \cmd{read\_saif} 读入 SAIF，或用 \cmd{set\_switching\_activity} 在 net 上注解开关活动（switching activity）。更多说明见 Annotating the Switching Activity。
  \item 将应用选项 \appopt{abstract.annotate\_power} 设为 \texttt{true}，以便在 Abstract 中存储功耗信息。
  \item 用 \cmd{create\_abstract} 为该 block 创建 Abstract。
\end{enumerate}

创建带功耗信息的 Abstract 时，工具只为被移除的逻辑存储功耗信息。

在顶层例化此类 Abstract 并做功耗分析时，工具会：
\begin{itemize}
  \item 根据顶层看到的开关活动与负载，在上下文中重新计算 Abstract 内接口逻辑的功耗；
  \item 对被 Abstract 创建时移除的逻辑，使用 Abstract 中已存储的功耗。
\end{itemize}

要用 \cmd{report\_power -blocks} 报告顶层例化 Abstract 中的功耗信息。若存在多层物理层次，可用 \opt{-levels} 指定要报告的层数。

\subsection{为信号电迁移分析创建 Abstract}
\subsubsection*{Creating Abstracts for Signal Electromigration Analysis}

要创建可用于顶层信号电迁移（signal electromigration）分析的 Abstract，在运行 \cmd{create\_abstract} 之前，将应用选项 \appopt{abstract.enable\_signal\_em\_analysis} 设为 \texttt{true}。

\subsection{处理多层物理层次}
\subsubsection*{Handling Multiple Levels of Physical Hierarchy}

对于多层物理层次设计，若要为内部已例化 Abstract 的 block 再创建 Abstract，须先：
\begin{itemize}
  \item 将希望用 Abstract 表示的下层 block 绑定到特定 Abstract。
\end{itemize}

例如，假设设计具有如下多层物理层次：

\begin{figure}[htbp]
\centering
\fbox{\parbox{0.85\textwidth}{\centering\vspace{1.5cm}
\textit{（原书 Figure 192：A Design With Multiple Levels of Physical Hierarchy）}\\[0.3em]
\small BOT / MID / TOP 层次关系示意，请参阅原版 PDF 插图。
\vspace{1.5cm}}}
\caption{多层物理层次设计}
\label{fig:hier-multi-level}
\end{figure}

要为名为 \texttt{BOT} 的 block 创建 Abstract：
\begin{lstlisting}
fc_shell> open_block BOT
fc_shell> create_abstract
\end{lstlisting}

要为名为 \texttt{MID} 的 block 创建 Abstract：
\begin{lstlisting}
fc_shell> open_block MID
fc_shell> change_abstract -view abstract -references BOT
fc_shell> create_abstract
\end{lstlisting}

默认情况下，\cmd{create\_abstract} 保留所有物理层次。要用 \opt{-preserve\_block\_instances false} 展平物理层次并仅保留下层 block 的相关逻辑，例如：
\begin{lstlisting}
fc_shell> open_block MID
fc_shell> create_abstract -preserve_block_instances false
\end{lstlisting}

在下层 block 按只读处理的顶层实现流程中使用展平 Abstract。展平 Abstract 可减小顶层网表规模，并简化数据管理——无需把下层 block 放进顶层的参考库列表。

\subsection{为最小脉宽分析创建 Abstract}
\subsubsection*{Creating Abstracts for Minimum Pulse Width Analysis}

在顶层时序 Abstract 中，Fusion Compiler 支持检查最小脉宽（minimum pulse width，MPW）违规。创建 Abstract 时，工具会保留存在 MPW 违规的时钟 sink，及其扇入 cell 与 net。

可用下列应用选项：
\begin{itemize}
  \item 将 \appopt{abstract.keep\_min\_pulse\_width\_violation\_pins} 设为 \texttt{true}，检查并保留带 MPW 违规的寄存器及其完整时钟路径。这些寄存器的数据路径不会保留。该选项默认值为 \texttt{false}。
  \item 将 \appopt{abstract.min\_pulse\_width\_violation\_slack\_limit} 设为所需 slack，以限制纳入的 MPW 违规范围。这有助于保留尚未违规、但 slack 略为正的 MPW pin。
\end{itemize}

% ============================================================
\section{报告 Abstract 纳入原因}
\subsection*{Reporting Abstract Inclusion Reasons}

在 Fusion Compiler 中，\cmd{create\_abstract} 会在顶层与块级 Abstract 视图中纳入 cell、net、pin 等网表对象，并报告纳入原因。

要用 \cmd{report\_abstract\_inclusion\_reason} 报告特定网表对象被纳入 Abstract 视图的原因；该命令可覆盖 Abstract 中各层次的对象。

每条原因会显示一个原因码（reason code），对应有效的纳入理由。同一网表对象可有多个原因。

\begin{noteBox}
\cmd{report\_abstract\_inclusion\_reason} 不会报告层次 cell，或层次 cell 上 pin 的纳入原因。
\end{noteBox}

可在例化了 Abstract 的顶层设计上做原因报告，无需把 Abstract 设为当前 block。

下例报告 Abstract 视图中 pin 的纳入原因：
\begin{lstlisting}
fc_shell> report_abstract_inclusion_reason [get_pins -of_object \
 [get_cells -filter "is_hierarchical==false" -hierarchical *]]

Legend
mv - MV Logic
comp - Compact interface logic
cell - Cell inclusion
--------------------------------------------------------------------
Pin name                            Reason(s) for inclusion
--------------------------------------------------------------------
AINV_P_260/I                        comp, cell
AINV_P_260/VDD                      cell
AINV_P_260/VSS                      cell
AINV_P_260/ZN                       comp, cell
AINV_P_262/I                        comp, cell
AINV_P_262/VDD                      cell
AINV_P_262/VSS                      cell
AINV_P_262/ZN                       comp, cell
AINV_P_263/I                        mv, comp, cell
\end{lstlisting}

报告开头会显示图例（legend），说明各原因码的详细含义。

% ============================================================
\section{创建 Abstract 之后修改 Block}
\subsection*{Making Changes to a Block After Creating an Abstract}

用 \cmd{create\_abstract} 为 block 创建 Abstract 视图时，默认会同时保存该 block 的 Abstract 与 design 视图，并把 design 视图改为只读。

可以打开只读 block 并在内存中修改，但不能用 \cmd{save\_block} 或 \cmd{save\_lib} 保存这些修改。要保存对只读 block 的修改，可用下列方法之一：
\begin{itemize}
  \item 用 \cmd{save\_block -as} 另存为不同的 block；
  \item 用 \cmd{create\_abstract} 为该 block 重新创建 Abstract，工具会自动保存 design 视图；
  \item 若不再需要 Abstract，用 \cmd{remove\_abstract} 删除它。这会把 design 视图从只读改回可编辑，随后可用 \cmd{save\_block} 或 \cmd{save\_lib} 保存。
\end{itemize}

若用 \cmd{create\_abstract -read\_only} 创建只读 Abstract，工具不会把该 block 的 design 视图设为只读，后续修改可用 \cmd{save\_block} 或 \cmd{save\_lib} 保存，例如：
\begin{lstlisting}
fc_shell> create_abstract -read_only
fc_shell> change_link [get_cells U25] AND2
fc_shell> save_block
\end{lstlisting}

% ============================================================
\section{创建 Frame 视图}
\subsection*{Creating a Frame View}

要执行顶层布线（含 virtual routing），每个 Abstract 视图都必须有对应的 Frame 视图。用 \cmd{create\_frame} 创建 Frame 视图；该命令从 design 视图提取 blockage、pin 与 via 信息。

创建 Frame 视图后，用 \cmd{save\_lib} 保存设计库。

% ============================================================
\section{在顶层链接 Abstract 视图}
\subsection*{Linking to Abstract Views at the Top-Level}

要使顶层设计链接到下层 block 的 Abstract 视图，包含该 Abstract 视图的设计库必须是顶层的参考库之一。

例如，假设顶层设计名为 \texttt{TOP}，并例化了两个下层 block \texttt{BLK1} 与 \texttt{BLK2}，如图~\ref{fig:hier-top-blocks} 所示。

\begin{figure}[htbp]
\centering
\fbox{\parbox{0.85\textwidth}{\centering\vspace{1.5cm}
\textit{（原书 Figure 193：Top-Level Design With Instantiated Blocks）}\\[0.3em]
\small TOP 例化 BLK1、BLK2 示意，请参阅原版 PDF 插图。
\vspace{1.5cm}}}
\caption{例化了 Block 的顶层设计}
\label{fig:hier-top-blocks}
\end{figure}

要链接到 \texttt{BLK1} 与 \texttt{BLK2} 的 Abstract，包含这些 Abstract 视图的设计库必须是 \texttt{TOP} 的参考库。若尚未加入，可用 \cmd{set\_ref\_libs -add} 添加，例如：
\begin{lstlisting}
current_block TOP
set_ref_libs -add ../BLK1/BLK1.nib
set_ref_libs -add ../BLK2/BLK2.nib
\end{lstlisting}

要用 \cmd{report\_ref\_libs} 报告当前设计的参考库。

链接顶层设计时，若下层 block 已有 Abstract 视图，工具默认链接到该 Abstract。链接下层 block 时各视图的默认优先级为：
\begin{enumerate}
  \item Abstract 视图
  \item Design 视图
  \item Frame 视图
  \item Outline 视图
\end{enumerate}

% ============================================================
\section{将 Block 从 Abstract 切换为 Design 视图}
\subsection*{Changing the Block From Abstract to Design View}

要用 \cmd{change\_abstract}，并通过 \opt{-view} 指定目标视图，把 block 从 Abstract 切换为 design 视图。从 Abstract 切到 design 视图时，应移除现有约束并施加全芯片时序约束。

下例将 \texttt{BLK1} 从 Abstract 切到 design 视图，并施加全芯片时序约束。
\begin{enumerate}
  \item 移除现有时序约束：
\begin{lstlisting}
fc_shell> remove_scenarios -all
fc_shell> remove_modes -all
fc_shell> remove_corners -all
\end{lstlisting}
  \item 用 \cmd{change\_abstract} 从 Abstract 切到 design 视图：
\begin{lstlisting}
fc_shell> change_abstract -view design -references BLK1
\end{lstlisting}
  \item 执行全芯片 scenario 创建脚本：
\begin{lstlisting}
fc_shell> source full_chip_scenario_creation.tcl
\end{lstlisting}
\end{enumerate}

若 block 在其设计库中的 Abstract 视图已更新，可用 \cmd{change\_abstract -reload} 重新加载新的 Abstract。

要用 \cmd{report\_abstracts} 报告设计当前链接到的 Abstract 视图。

关于指定参考库以及在顶层链接 Abstract 视图的更多信息，见 SolvNet 文章 2436119，\textit{Setting Up Designs for Hierarchical Place and Route}。

% ============================================================
\section{链接带多个 Label 的子块}
\subsection*{Linking to Subblocks With Multiple Labels}

在实现流程的不同阶段保存 block 时，可用 \cmd{save\_block} 的 \opt{-as blockName/labelName} 或 \opt{-label labelName}，为每次保存指定唯一 label。

默认情况下，链接时顶层会链接到具有相同 label 的下层 block。若子块已用多个 label 保存，可用 \cmd{set\_label\_switch\_list} 指定顶层应链接到哪个 label。该命令给出链接时按优先级排序的 label 列表。默认作用于所有子块；要用 \opt{-reference} 指定不同的 block。

下例指定顶层应链接到 \texttt{BLOCK1} 与 \texttt{BLOCK2} 的 \texttt{PostCompile} label：
\begin{lstlisting}
fc_shell> set_label_switch_list \
    -reference {BLOCK1 BLOCK2} PostCompile
\end{lstlisting}

下例指定：在顶层链接任意子块时，\texttt{PostCompile} 优先于 \texttt{PreCompile}：
\begin{lstlisting}
fc_shell> set_label_switch_list {PostCompile PreCompile}
\end{lstlisting}

可再次使用 \cmd{set\_label\_switch\_list} 更新 block 的 label 切换列表，再用 \cmd{link\_block -rebind} 重新链接。重新链接时注意：
\begin{itemize}
  \item 若切换列表中的 label 在某子块均不可用，或指定了空切换列表，工具会把该子块链接到其先前已链接的 label；
  \item 若已移除某子块的全部 label，工具会把该子块链接到其父 block 的 label。
\end{itemize}

% ============================================================
\section{从顶层指定 Block 的可编辑性}
\subsection*{Specifying the Editability of Blocks From the Top-Level}

对于带物理层次的设计，可在顶层用 \cmd{set\_editability} 指定下层 block 是否可修改。

可更改：
\begin{itemize}
  \item 用 \opt{-blocks} 指定的特定 block；
  \item 用 \opt{-from\_level} 指定从某一物理层次起的全部 block；
  \item 用 \opt{-to\_level} 指定到某一物理层次为止的全部 block。
\end{itemize}

对于使用只读 Abstract 的顶层实现流程，链接设计后应显式把下层 block 设为只读：
\begin{lstlisting}
fc_shell> set_editability -blocks [get_blocks -hierarchical] \
   -value false
\end{lstlisting}

默认情况下，若用下列命令之一更改库名或 block 名，工具不会把可编辑性设置从旧参考传播到新参考：
\begin{itemize}
  \item \cmd{change\_abstract -lib}
  \item \cmd{link\_block -rebind -force}
  \item \cmd{set\_reference}
\end{itemize}

要保留可编辑性设置，在更改库名或 block 名之前，将应用选项 \appopt{design.preserve\_reference\_editability} 设为 \texttt{true}。

用 \cmd{save\_block} 保存层次化设计时，可编辑性设置会被保存。

% ============================================================
\section{使用 Abstract 做顶层收敛前的准备}
\subsection*{Preparing for Top-Level Closure With Abstracts}

在顶层执行综合、布局、优化、时钟树综合与布线之前，必须完成下列任务：
\begin{itemize}
  \item 施加顶层时序约束与设置。\\
可用 \cmd{split\_constraints} 把芯片级约束拆成独立的顶层与块级约束。

  \item 用 \cmd{set\_block\_to\_top\_map} 审查并指定顶层与块级 mode、corner、clock 之间的对应关系。

  \item 若用于顶层布局与优化的是已经过 CCD 优化的预时钟树实现 block，则 CCD、clock-gate estimation 等引擎在 block 内寄存器或集成时钟门控（integrated clock gating，ICG）单元时钟 pin 上动态调整的理想时钟延迟（ideal clock latency），在顶层可能看不到，从而导致顶层时序不准确。要用下列命令把这些时钟延迟约束从块级提升到顶层：
\begin{lstlisting}
promote_clock_data -latency_offset -auto_clock connected
\end{lstlisting}

  \item 施加顶层时钟树综合设置与例外：
  \begin{itemize}
    \item 若顶层例外未包含下层 block 内 pin 上的 balance point，用下列命令从下层提升 balance point：
\begin{lstlisting}
promote_clock_data -auto_clock connected -balance_points
\end{lstlisting}
    \item 若在块级执行了并发时钟与数据（concurrent clock and data，CCD）优化，工具会为块级时钟端口推导中位时钟延迟（median clock latency），以反映 CCD 期间为块级寄存器推导的延迟调整。要用下列命令把这些块级时钟延迟作为块级时钟 pin 的 clock balance point delay 提升：
\begin{lstlisting}
promote_clock_data -auto_clock connected -port_latency
\end{lstlisting}
    \item 若下层 block 含 clock mesh，用 \cmd{promote\_clock\_data -mesh\_annotations} 提升 mesh 注解（已注解的 transition 与 delay）。
  \end{itemize}

  \item 若只需为特定顶层时钟提升时钟数据，使用：
\begin{lstlisting}
promote_clock_data -clocks {top_clock} -port_latency
\end{lstlisting}
该命令也会提升与时钟无关的例外。若给定顶层时钟在任一 mode 中都未映射到任何块级时钟，工具会发出警告。

  \item 若顶层存在无法映射到块级对应 mode/corner/clock 的活动 mode、corner 或 clock，仍可用下列命令为其他已映射的 mode/corner/clock 提升时钟数据：
\begin{lstlisting}
promote_clock_data -balance_points -force
\end{lstlisting}
该命令会对未映射的 mode、corner 或 clock 发出错误。

  \item 施加顶层 UPF 约束；可用 \cmd{split\_constraints} 从全芯片 UPF 约束生成。\\
工具会把下层 block 中所需的 UPF 约束提升到顶层。
\end{itemize}

% ============================================================
\section{检查含 Abstract 设计的顶层收敛问题}
\subsection*{Checking Designs With Abstracts for Top-Level-Closure Issues}

可用 \cmd{check\_hier\_design} 检查含 Abstract 的层次化设计，查找应用选项一致性问题，以及与顶层收敛相关的设计问题。在执行顶层收敛前识别并修复这些问题，有助于缩短周转时间（turnaround time）。

使用 \cmd{check\_hier\_design} 时可指定：
\begin{itemize}
  \item 用 \opt{-reference} 指定要检查的参考。若不指定，命令检查物理层次中的所有参考。
  \item 用 \opt{-stage} 指定检查类型：
  \begin{itemize}
    \item \opt{-stage timing}：执行时序相关检查。若不指定 \opt{-stage}，默认执行时序相关检查。
    \item \opt{-stage pre\_placement}：同时执行时序与预布局相关检查。
  \end{itemize}
\end{itemize}

\cmd{check\_hier\_design} 可检查顶层约束与每个链接到 Abstract 的下层实例约束之间的一致性。启用该功能时，在运行 \cmd{check\_hier\_design} 之前将应用选项 \appopt{abstract.check\_constraints\_consistency} 设为 \texttt{true}。

工具在 \cmd{create\_abstract} 期间会保存预定列表中的 timer 应用选项设置。在顶层，\cmd{check\_hier\_design} 会将这些设置与块级设计中的对应设置进行比较。

下列应用选项会在 \cmd{create\_abstract} 期间自动保存，并由 \cmd{check\_hier\_design} 校验：
\begin{itemize}
  \item \appopt{time.case\_analysis\_propagate\_through\_icg}
  \item \appopt{time.case\_analysis\_sequential\_propagation}
  \item \appopt{time.clock\_gating\_propagate\_enable}
  \item \appopt{time.clock\_gating\_user\_setting\_only}
  \item \appopt{time.clock\_marking}
  \item \appopt{time.clock\_reconvergence\_pessimism}
  \item \appopt{time.create\_clock\_no\_input\_delay}
  \item \appopt{time.crpr\_remove\_clock\_to\_data\_crp}
  \item \appopt{time.delay\_calc\_waveform\_analysis\_mode}
  \item \appopt{time.delay\_calculation\_style}
  \item \appopt{time.disable\_case\_analysis\_ti\_hi\_lo}
  \item \appopt{time.disable\_clock\_gating\_checks}
  \item \appopt{time.disable\_cond\_default\_arcs}
  \item \appopt{time.disable\_internal\_inout\_net\_arcs}
  \item \appopt{time.disable\_recovery\_removal\_checks}
  \item \appopt{time.edge\_specific\_source\_latency}
  \item \appopt{time.enable\_auto\_mux\_clock\_exclusivity}
  \item \appopt{time.enable\_ccs\_rcv\_cap}
  \item \appopt{time.enable\_clock\_propagation\_through\_preset\_clear}
  \item \appopt{time.enable\_clock\_propagation\_through\_three\_state\_enable\_pins}
  \item \appopt{time.enable\_clock\_to\_data\_analysis}
  \item \appopt{time.enable\_non\_sequential\_checks}
  \item \appopt{time.enable\_preset\_clear\_arcs}
  \item \appopt{time.enable\_si\_timing\_windows}
  \item \appopt{time.gclock\_source\_network\_num\_master\_registers}
  \item \appopt{time.special\_path\_group\_precedence}
  \item \appopt{time.use\_lib\_cell\_generated\_clock\_name}
  \item \appopt{time.use\_special\_default\_path\_groups}
\end{itemize}

关于 \cmd{check\_hier\_design} 识别的问题及修复方法，请参阅对应消息 ID 的 man page。

除生成报告外，\cmd{check\_hier\_design} 还会生成增强消息系统（enhanced messaging system，EMS）数据库，可在 Fusion Compiler GUI 的消息浏览器中查看。应在运行 \cmd{check\_hier\_design} 之前创建 EMS 数据库，例如：
\begin{lstlisting}
fc_shell> create_ems_database check_hier.ems
fc_shell> check_hier_design -stage timing
fc_shell> save_ems_database
\end{lstlisting}

在 Fusion Compiler GUI 消息浏览器中，可对消息排序、过滤，并链接到对应 man page，如图~\ref{fig:hier-ems-browser} 所示。

\begin{figure}[htbp]
\centering
\fbox{\parbox{0.85\textwidth}{\centering\vspace{1.5cm}
\textit{（原书 Figure 194：Viewing the EMS Database in the Message Browser）}\\[0.3em]
\small 请参阅原版 PDF 插图。
\vspace{1.5cm}}}
\caption{在消息浏览器中查看 EMS 数据库}
\label{fig:hier-ems-browser}
\end{figure}

也可用 \cmd{report\_ems\_database} 以 ASCII 格式输出 EMS 数据库中的信息。

\cmd{check\_hier\_design} 执行的检查同样可在 \cmd{check\_design} 中使用。可以：
\begin{itemize}
  \item 用 \cmd{check\_design -checks hier\_timing} 执行专用于顶层收敛的时序检查；
  \item 用 \cmd{check\_design -checks timing} 执行全部时序检查（含顶层收敛专用检查）；
  \item 用 \cmd{check\_design -checks hier\_preplacement} 执行专用于顶层收敛的预布局检查；
  \item 用 \cmd{check\_design -checks preplacement} 执行全部预布局检查（含顶层收敛专用检查）。
\end{itemize}

也可为 \cmd{check\_design} 生成 EMS 数据库并在 GUI 中查看，方式与 \cmd{check\_hier\_design} 类似。

\subsection{用 Early Data Check Manager 处理设计数据}
\subsubsection*{Handling Design Data Using the Early Data Check Manager}

在识别并修复顶层收敛相关问题的同时，工具允许在存在设计数据违规的情况下探索顶层流程。层次化检查可检测设计中各类不匹配、不完整与不一致的数据。可为层次化检查配置下列策略（policy）以管理设计数据违规：
\begin{itemize}
  \item \textbf{Error}：不缓解违规，而是发出错误消息。
  \item \textbf{Tolerate}：做出可预期、可解释的假设以继续流程；不对设计或设置做更改。
  \item \textbf{Repair}：用一种或多种修复策略缓解违规并记录。修复可包括对设计与设置的更改。
  \item \textbf{Strict}：按检查所支持的策略，依次尝试 error、tolerate、repair。
  \item \textbf{Lenient}：按检查所支持的策略，依次尝试 repair、tolerate、error。
\end{itemize}

\begin{noteBox}
对错误的容忍（toleration）与修复（repair）可能影响顶层流程的 QoR。错误容忍与修复主要用于早期设计探索阶段管理数据违规。在实现流程中，应对所有检查使用 strict 策略。
\end{noteBox}

\subsubsection{处理早期设计数据的前提}
\subsubsection*{Prerequisites for Handling Early Design Data}

工具在 \cmd{check\_hier\_design} 命令内部执行修复，并把记录保存到子块中。对于需要重建 Frame 的修复，还需运行 \cmd{save\_lib} 把 Frame 保存到磁盘。

要将修复及其记录成功保存到子块，需要下列设置：
\begin{enumerate}
  \item 以可编辑方式打开包含子块的参考库：
\begin{lstlisting}
open_lib top.nlib -ref_libs_for_edit
\end{lstlisting}
  \item 对需要执行修复的子块启用可编辑性：
\begin{lstlisting}
set_editability -blocks <> -value true
\end{lstlisting}
\end{enumerate}

若需执行修复，\cmd{check\_hier\_design} 会检查所需可编辑性设置；若未正确设置，会发出 \texttt{TL-170} 错误消息。

\begin{noteBox}
工具支持对链接到 \texttt{read\_only} Abstract（而非可编辑 Abstract）的子块执行修复，因为管理不匹配设计数据的机制目标只是让顶层探索流程能继续，子块 Abstract 上的修复工作不应合并回实际 design 视图。若顶层使用可编辑 design 视图，修复工作仍会继续。
\end{noteBox}

\subsubsection{早期数据检查、策略与修复策略}
\subsubsection*{Early Data Checks, Policies, and Strategies}

工具可在设计早期阶段识别数据问题与违规。可为预定义检查设置策略与策略配置，以允许、容忍或修复数据。表~\ref{tab:hier-tlc-checks} 描述了顶层收敛的预定义检查。

\begin{table}[htbp]
\centering
\caption{顶层收敛检查、策略与修复策略（原书 Table 64）}
\label{tab:hier-tlc-checks}
\small
\begin{tabularx}{\textwidth}{@{}>{\ttfamily\raggedright\arraybackslash}p{0.28\textwidth} >{\raggedright\arraybackslash}X >{\raggedright\arraybackslash}p{0.22\textwidth} >{\raggedright\arraybackslash}p{0.12\textwidth} >{\raggedright\arraybackslash}p{0.1\textwidth}@{}}
\toprule
Check & 说明 & Strategy & 支持 Policies & 支持 References \\
\midrule
hier.block.missing\_frame\_view
  & 检查层次 block 是否缺少 Frame 视图
  & 从已绑定视图重建 Frame
  & error, tolerate, repair & MID, BOT \\
hier.block.reference\_missing\_port\_location
  & 检查物理层次边界 pin 位置是否缺失
  & 按连接关系在 block 边界上为 block pin 分配位置
  & error, tolerate, repair & MID, BOT \\
hier.block.reference\_port\_outside\_boundary
  & 检查物理层次边界 pin 位置是否在边界外
  & 按连接关系在 block 边界上重新分配位置
  & error, tolerate, repair & MID, BOT \\
hier.block.reference\_missing\_port
  & 检查物理层次参考中是否缺少 port
  & 在参考 block 中创建缺失 port 并更新位置
  & error, tolerate, repair & MID, BOT \\
hier.block.instance\_bound\_to\_frame
  & 检查物理层次实例是否绑定到 Frame 视图
  & --- & error, tolerate & TOP \\
hier.block.instance\_with\_design\_type\_macro
  & 检查物理层次实例是否绑定到 macro
  & --- & error, tolerate & TOP \\
hier.top.estimated\_corner
  & 检查顶层是否存在 estimated corner
  & --- & error, tolerate & TOP \\
hier.block.missing\_leaf\_cell\_location
  & 检查物理层次 leaf cell 位置是否缺失
  & 在 block 边界内为 leaf cell 分配近似位置（不考虑合法性）
  & error, tolerate, repair & MID, BOT \\
hier.block.leaf\_cell\_outside\_boundary
  & 检查物理层次 leaf cell 位置是否在边界外
  & 将 leaf cell 重新放到边界内最近位置（不考虑合法性）
  & error, tolerate, repair & MID, BOT \\
hier.block.missing\_child\_block\_instance\_location
  & 检查指定物理层次中子物理层次实例位置是否缺失
  & 按连接关系在 block 边界内为下层子块分配近似位置（不考虑重叠消除）
  & error, tolerate, repair & MID \\
hier.block.child\_block\_instance\_outside\_boundary
  & 检查子物理层次实例位置是否在指定物理层次边界外
  & 将子块移到边界内最近位置（不考虑重叠消除）
  & error, tolerate, repair & MID \\
hier.block.port\_mismatch\_between\_views
  & 检查物理层次各视图之间是否存在 port 不匹配
  & 以绑定视图为 golden，重建 Frame 以与之匹配
  & error, tolerate, repair & MID, BOT \\
hier.block.missing\_design\_view\_for\_abs
  & 检查 Abstract 是否缺少 design 视图
  & --- & error, tolerate & MID, BOT \\
hier.block.instance\_unlinked
  & 检查物理层次参考是否因视图从 Frame/ETM 变为 Abstract/design 而未正确链接到顶层
  & --- & error & TOP \\
hier.block.abstract\_type\_non\_timing
  & 检查 Abstract 类型是否非 timing
  & --- & error & MID, BOT \\
hier.block.abstract\_target\_use\_non\_implementation
  & 检查 Abstract 的 target use 是否非 implementation
  & --- & error & MID, BOT \\
hier.block.missing\_core\_area
  & 检查物理层次是否缺少 core area
  & --- & error & MID, BOT \\
hier.block.unmapped\_logic
  & 检查物理层次中是否存在未映射逻辑
  & --- & error & MID, BOT \\
\bottomrule
\end{tabularx}
\end{table}

其中：
\begin{itemize}
  \item BOT 表示底层设计；
  \item MID 表示例化 BOT 的中间层设计；
  \item TOP 表示例化 MID 的顶层设计。
\end{itemize}

\subsubsection{为早期数据检查设置策略}
\subsubsection*{Setting the Policy for Early Data Checks}

要用 \cmd{set\_early\_data\_check\_policy} 设置或修改全部检查的策略，或修改数据检查的策略设置。可用下列选项配置策略：
\begin{itemize}
  \item \opt{-checks}：指定预定义数据检查列表。检查名支持星号通配符（\texttt{*}）。
  \item \opt{-policy}：按检查所支持的类型指定策略，可为：
  \begin{itemize}
    \item \texttt{error}：若检查不支持 \texttt{error} 策略，工具发出错误消息。
    \item \texttt{tolerate}：若检查不支持 \texttt{tolerate} 策略，工具发出错误消息。
    \item \texttt{repair}：若检查不支持 \texttt{repair} 策略，工具发出错误消息。
    \item \texttt{strict}：按下列顺序设置第一个匹配的策略：1.~error；2.~tolerate；3.~repair。例如，若检查支持 \texttt{error} 与 \texttt{tolerate}，选择 \texttt{strict} 时会设为 \texttt{error}。
    \item \texttt{lenient}：按下列顺序设置第一个匹配的策略：1.~repair；2.~tolerate；3.~error。例如，若检查支持 \texttt{error} 与 \texttt{tolerate}，选择 \texttt{lenient} 时会设为 \texttt{tolerate}。
  \end{itemize}
\begin{noteBox}
若同时指定 checks 与 policies，配置支持本节所述、且该应用特定检查所支持的全部策略。但若只指定 policies 而不指定 checks，则配置仅支持 \texttt{strict} 与 \texttt{lenient} 策略。
\end{noteBox}
  \item 也可在 \opt{-policy} 中使用 \texttt{if\_not\_exists} 关键字，仅当当前设计尚未设置策略时才设置。
  \item \opt{-strategy}：为检查指定要应用的 strategy。
\begin{noteBox}
\opt{-strategy} 不适用于层次化检查。
\end{noteBox}
  \item \opt{-references}：为参考子设计指定策略与 strategy。
\end{itemize}

\subsubsection{报告早期数据检查记录}
\subsubsection*{Reporting Early Data Check Records}

可用 \cmd{report\_early\_data\_checks} 报告全部检查或特定检查集的策略设置。可用下列选项配置报告：
\begin{itemize}
  \item \opt{-policy}：报告全部检查的配置详情，显示各检查所选策略与 strategy。与 \opt{-hierarchical} 合用可报告子块配置。
  \item \opt{-checks}：报告各检查的当前策略、支持的策略、strategy 与帮助文本。检查名支持星号通配符（\texttt{*}）。与 \opt{-hierarchical} 合用可报告子块检查。
  \item \opt{-verbose}：报告所有失败检查的摘要、策略、已应用 strategy，以及带注释的已检查对象。与 \opt{-hierarchical} 合用可报告子块中失败检查的记录、策略、strategy 与已检查对象。
  \item \opt{-hierarchical}：遍历层次，报告层次中每个设计的检查、策略、strategy 与失败计数。
\end{itemize}

下例中，\cmd{report\_early\_data\_checks -hierarchical -policy} 报告顶层与块级设计上层次化检查的策略设置：
\begin{lstlisting}
fc_shell> report_early_data_checks -policy
****************************************
Report : report_early_data_checks
Design : top
Version: R-2020.09
Date   : Thu Aug 20 23:53:54 2020
****************************************
Design                   Check                                             Policy
 Strategy
---------------------------------------------------------------------------------
-------
unit_des.nlib:top.design hier.block.instance_bound_to_frame                error

unit_des.nlib:top.design hier.block.instance_unlinked                      error

unit_des.nlib:top.design hier.block.instance_with_design_type_macro        error

unit_des.nlib:top.design hier.top.estimated_corner                         error
---------------------------------------------------------------------------------
------
blk                      hier.block.abstract_missing_design_view           error

blk                      hier.block.abstract_target_use_non_implementation error

blk                      hier.block.abstract_type_non_timing               error

blk                      hier.block.leaf_cell_outside_boundary             error
---------------------------------------------------------------------------------
-----
\end{lstlisting}

下例中，\cmd{report\_early\_data\_checks -hierarchical -verbose} 报告对已检查对象执行的修复，并给出子块上的策略、strategy 与已检查对象详情：
\begin{lstlisting}
fc_shell> report_early_data_checks -hierarchical -verbose
******************************************************
Design Check                               Policy Strategy Checked  Comment
                                                           Object
------------------------------------------------------------------------------------
blk  hier.block.missing_leaf_cell_location repair          eco_cell Instance location
                                                                    updated to
                                                                    X:150145 Y:69620
-------------------------------------------------------------------------------------
\end{lstlisting}

在运行顶层早期流程时，可能因子块或顶层设计中的设计数据违规而遇到顶层错误或警告。要利用数据检查与策略能力探索早期流程，需要知道设计上所见顶层错误对应的检查名。\cmd{report\_hier\_check\_description} 提供各顶层错误关联检查的信息。下例显示该命令生成的报告：
\begin{lstlisting}
fc_shell> report_hier_check_description
*********************************************
Report : Hier check description
Design : top
Version: R-2020.09
Date   : Fri Aug 21 00:08:32 2020
*********************************************

Legend
    E - error
    R - repair
    T - tolerate
    TOP - Top-Design
    BLK - Respective Block-Ref

--------------------------------------------------------------------------------------
---------------------------
Check                                      Error  Tolerate Repair Allowed  Allowed
 Description
                                            ID      ID      ID    Policies References
--------------------------------------------------------------------------------------
---------------------------
hier.block.missing_frame_view              TL-101  TL-401  TL-501  E|R|T     BLK
 Missing frame view
hier.block.abstract_missing_design_view    TL-101  TL-401  N/A     E|T       BLK
 Missing design view for abstract
hier.block.reference_missing_port_location TL-126  TL-426  TL-526  E|R|T     BLK
 Missing location of physical
                                                      hierarchy boundary pin
hier.block.reference_port_outside_boundary TL-127  TL-427  TL-527  E|R|T     BLK
 Location of physical hierarchy boundary
                                                      pin outside physical hierarchy boundary
\end{lstlisting}

\subsubsection{生成早期数据检查记录报告}
\subsubsection*{Generating a Report of Early Data Check Records}

可用 \cmd{get\_early\_data\_check\_records} 从设计生成检查记录报告。与 \opt{-hierarchical} 合用可获取所有子块的检查记录。

报告格式为：
\begin{lstlisting}
check_name@object
\end{lstlisting}

例如：
\begin{lstlisting}
fc_shell> get_early_data_check_records -hierarchical
{hier.block.missing_leaf_cell_location@eco_cell}
\end{lstlisting}

可用 \opt{-filter} 按对象类、检查名等参数过滤结果，例如：
\begin{lstlisting}
fc_shell> get_early_data_check_records \
-filter "checked_object_class==port"
{place.port_type_mismatch@clk route.missing_layer_direction@feed_in
 route.layer_name_mismatch@out}
\end{lstlisting}

% ============================================================
\section{使用 Abstract 视图执行顶层收敛}
\subsection*{Performing Top-Level Closure With Abstract Views}

可先实现各 block，再实现顶层，如图~\ref{fig:hier-tlc-after-blocks} 所示：

\begin{figure}[htbp]
\centering
\fbox{\parbox{0.85\textwidth}{\centering\vspace{1.5cm}
\textit{（原书 Figure 195：Top-Level Implemented After Blocks are Completed）}\\[0.3em]
\small 块级流程完成后生成 Abstract/Frame，再进入顶层综合/布局/优化/CTS/布线/后布线优化。请参阅原版 PDF 插图。
\vspace{1.5cm}}}
\caption{Block 完成后再实现顶层}
\label{fig:hier-tlc-after-blocks}
\end{figure}

也可并行实现 block 与顶层，如图~\ref{fig:hier-tlc-parallel} 所示：

\begin{figure}[htbp]
\centering
\fbox{\parbox{0.85\textwidth}{\centering\vspace{1.5cm}
\textit{（原书 Figure 196：Top and Block Levels Implemented in Parallel）}\\[0.3em]
\small 块级与顶层流程并行，中途交换 Abstract/Frame。请参阅原版 PDF 插图。
\vspace{1.5cm}}}
\caption{顶层与块级并行实现}
\label{fig:hier-tlc-parallel}
\end{figure}

在顶层使用 Abstract 时，可用块级支持的命令执行顶层综合、布局、优化、时钟树综合、布线与后布线优化。当前工具在顶层收敛期间不会修改 Abstract 内部，因此可创建并使用只读 Abstract。

% ============================================================
\section{创建 ETM 与 ETM Cell Library}
\subsection*{Creating ETMs and ETM Cell Libraries}

要在 Fusion Compiler 中创建 ETM 与 ETM cell library，见「在 Fusion Compiler 工具中创建 ETM 与 ETM Cell Library」。

要分别在 PrimeTime 与 Library Manager 中创建 ETM 与 ETM cell library，见：
\begin{itemize}
  \item 在 PrimeTime 工具中创建 ETM
  \item 在 Library Manager 工具中创建 ETM Cell Library
\end{itemize}

\subsection{在 Fusion Compiler 工具中创建 ETM 与 ETM Cell Library}
\subsubsection*{Creating ETMs and ETM Cell Libraries in the Fusion Compiler Tool}

不必先在 PrimeTime 中用 \cmd{extract\_model} 为每个 mode/corner 组合创建 ETM，再在 Library Manager 中把 ETM 与对应物理信息合并成 cell library；可直接在 Fusion Compiler 中用 \cmd{extract\_model} 一步完成这两步。

下列 Fusion Compiler 脚本为设计的每个 mode 与 corner 创建 ETM，再与对应物理信息合并并创建相应 cell library：
\begin{lstlisting}
# Open the design, create a frame view, and set the PrimeTime options
open block.nlib:block.design
create_frame <options>        ;# extract_model needs Frame for library
 preparation
set_pt_options -pt_exec_path <> -work_dir ETM_work_dir \
               -post_link_script <tcl script with extract_model*
 variables> \

# To use StarRC for Parasitic extraction
set_app_options -name extract.starrc_mode -value true
set_starrc_options -config <starrc_config_file>

# Create the ETM and generate the ETM cell library
extract_model
\end{lstlisting}

默认情况下，Fusion Compiler 将生成的 ETM library 写出到：
\begin{itemize}
  \item 带 label 的 block：\texttt{[pwd]/ETM\_Lib\_work\_dir/block\_name/label\_name/}
  \item 不带 label 的 block：\texttt{[pwd]/ETM\_Lib\_work\_dir/block\_name/}
\end{itemize}

\subsection{在 PrimeTime 工具中创建 ETM}
\subsubsection*{Creating ETMs in the PrimeTime Tool}

可在 PrimeTime 中用 \cmd{extract\_model} 为设计创建 ETM。对于多 corner-多 mode 设计，必须为每个 scenario 施加相应 corner 与 mode 约束并分别创建 ETM。

下列 PrimeTime 脚本为 AMS\_BLK 设计的 S3 scenario（由 m1 mode 与 c3 corner 组成）创建 ETM：
\begin{lstlisting}
#Read in the design
read_verilog ./AMS_BLK.v
link

# Apply parasitics
read_parasitics ./AMS_BLK.spef

#For multivoltage designs, apply UPF data and settings
load_upf ./AMD_BLK.upf
set extract_model_include_upf_data true

# Apply the mode (m1) and corner (c3) constraints for the scenario (S3)
source m1_constraints.tcl
source c3_constraints.tcl

# Enable clock latencies for designs with synthesized clock trees
set extract_model_with_clock_latency_arcs true
set extract_model_clock_latency_arcs_include_all_registers false

# Create the ETM
extract_model -library_cell -format db -output AMS_BLK_m1_c3
\end{lstlisting}

关于在 PrimeTime 中创建 ETM 的更多信息，见 \textit{PrimeTime User Guide} 的 Extracted Timing Models 一章。

\subsection{在 Library Manager 工具中创建 ETM Cell Library}
\subsubsection*{Creating ETM Cell Libraries in the Library Manager Tool}

为每个 mode/corner 组合创建 ETM 后，用 Library Manager 把 ETM 与对应物理信息合并并创建 cell library。更多信息见 \textit{Library Manager User Guide}。

下列 Library Manager 脚本把 AMS\_BLK 设计各 mode/corner 的 ETM 与对应物理信息合并，并创建相应 cell library：
\begin{lstlisting}
# Create a library work space
create_workspace -flow etm_moded AMS_BLK

# Read the physical data (frame view)
read_ndm -views frame AMS_BLK.ndm

# Read the ETMs for every scenario
read_db -mode_label m1 AMS_BLK_m1_c1.db
read_db -mode_label m1 AMS_BLK_m1_c2.db
read_db -mode_label m1 AMS_BLK_m1_c3.db
read_db -mode_label m2 AMS_BLK_m2_c1.db
read_db -mode_label m2 AMS_BLK_m2_c2.db
read_db -mode_label m2 AMS_BLK_m2_c3.db


# Check the compatibility of the libraries you read in
check_workspace

# Generate cell library
commit_workspace -output AMS_BLK_ETM.ndm
\end{lstlisting}

创建 ETM 的 cell library 时，可用下列方法之一获取物理数据：
\begin{itemize}
  \item 用 \cmd{read\_ndm -views frame} 读入对应 block 的 Frame 视图；
  \item 用 \cmd{read\_lef} 读入 block 的 LEF 文件；
  \item 用 \cmd{read\_gds} 读入 block 的 GDSII 文件；
  \item 用 \cmd{read\_oasis} 读入 block 的 OASIS 文件。
\end{itemize}

% ============================================================
\section{在顶层链接 ETM}
\subsection*{Linking to ETMs at the Top Level}

要在顶层链接 ETM，在创建顶层设计库之前把 ETM library 加入参考库列表，例如：
\begin{lstlisting}
fc_shell> lappend nt_ref_lib "AMS_BLK_ETM.ndm"
fc_shell> create_lib -technology tech.tf -ref_libs $nt_ref_lib TOP
\end{lstlisting}

包含多个 mode 的 ETM 称为 moded ETM。对于链接到 moded ETM 的 cell 实例，必须用 \cmd{set\_cell\_mode} 指定所需 mode，例如：
\begin{lstlisting}
fc_shell> current_mode m1
fc_shell> set_cell_mode m1 U1
\end{lstlisting}

对于链接到 moded ETM 的 cell 实例，若不指定 mode，工具不会激活时序弧、时序检查、生成时钟与 case value。

要用 \cmd{report\_cell\_modes} 报告特定 cell 实例的 cell mode。

要用 \cmd{set\_reference} 把 cell 实例在其 Abstract 视图与 ETM 之间切换。下例把 \texttt{U1} 从 Abstract 切到 ETM：
\begin{lstlisting}
fc_shell> set_reference -block AMS_BLK_ETM.ndm:AMS_BLK.timing U1
\end{lstlisting}

下例把 \texttt{U1} 从 ETM 切回 Abstract：
\begin{lstlisting}
fc_shell> set_reference -block AMS_BLK.ndm:AMS_BLK.abstract U1
\end{lstlisting}

在 ETM 与 Abstract 视图之间切换后，需重新施加顶层时序约束。

% ============================================================
\section{使用 ETM 执行顶层收敛}
\subsection*{Performing Top-Level Closure With ETMs}

在顶层使用 ETM 之前，必须完成下列任务：
\begin{itemize}
  \item 施加仅顶层的时序约束与设置。\\
若由 ETM 表示的 cell 实例含内部时钟，必须从顶层重新施加这些时钟定义。避免引用 block 内部对象的跨边界时序例外。

  \item 要施加 UPF 约束与设置，必须：
  \begin{enumerate}
    \item 用 \cmd{save\_upf} 从各块级设计生成块级 UPF；
    \item 在顶层，为每个由 ETM 表示的 block 实例加载块级 UPF，并加载仅顶层 UPF：
\begin{lstlisting}
load_upf -scope <instance_name> block.upf
load_upf top.upf
\end{lstlisting}
  \end{enumerate}

  \item 施加顶层布局约束与设置，例如 hard keepout margin。
\end{itemize}

在顶层使用 ETM 时，可用块级支持的命令执行综合、顶层布局、优化、时钟树综合、布线与后布线优化。

% ============================================================
\section{透明层次优化}
\subsection*{Transparent Hierarchy Optimization}

透明层次优化（Transparent Hierarchy Optimization，THO）使工具能透明地优化层次化设计。

THO 允许在遵守物理 block 边界的同时，并发优化顶层与子块，并使用 CCD、AWP、PrimeTime 延迟计算等全部高级优化技术。

在 THO 期间，工具在顶层时序上下文中优化顶层与特定子块（design 视图）。THO 不改变子块形状与物理层次。默认情况下，工具遵守 pin 位置，但提供选项以重新分配 pin 位置以改善 QoR。

由于无需展平设计，可继续用层次化 block 做 signoff：signoff extraction 与 PrimeTime STA 均可层次化运行。此外，后期 ECO 也往往更容易实现。

THO 可用于 \cmd{compile\_fusion} 的 \texttt{final\_place} 与 \texttt{final\_opto} 阶段，以及 \cmd{route\_opt} 与 \cmd{hyper\_route\_opt}。

\begin{figure}[htbp]
\centering
\fbox{\parbox{0.85\textwidth}{\centering\vspace{1.5cm}
\textit{（原书 Figure 197：Transparent Hierarchy Optimization Flow）}\\[0.3em]
\small THO 流程：\texttt{open\_block -ref\_libs\_for\_edit} → 启用 THO 设置 → \texttt{compile\_fusion -from final\_place} / \texttt{route\_opt} 或 \texttt{hyper\_route\_opt}。请参阅原版 PDF 插图。
\vspace{1.5cm}}}
\caption{透明层次优化流程}
\label{fig:hier-tho-flow}
\end{figure}

下列示例流程展示 THO 在整体实现流程中的用法。本例中，THO 作为典型层次化优化流程之上的额外增量优化步骤运行。

THO 流程包括以下步骤：
\begin{itemize}
  \item 访问库（Accessing Libraries）
  \item 设置透明层次优化（Setting up Transparent Hierarchy Optimization）
  \item 在 \cmd{compile\_fusion} 的 \texttt{final\_place} 阶段执行 THO
  \item 在 \cmd{compile\_fusion} 的 \texttt{final\_opto} 阶段执行 THO
  \item \cmd{compile\_fusion} 之后重新预算（Rebudgeting）
  \item 在 \cmd{route\_opt} 或 \cmd{hyper\_route\_opt} 期间执行 THO
\end{itemize}

\subsection{访问库}
\subsubsection*{Accessing Libraries}

启用 THO 设置之前，必须确保：
\begin{itemize}
  \item 所有子块库以编辑模式打开。必须用 \cmd{open\_lib} 的 \opt{-ref\_libs\_for\_edit} 打开父库。
  \item block 的 design 视图可编辑。
\end{itemize}

\subsection{设置透明层次优化}
\subsubsection*{Setting up Transparent Hierarchy Optimization}

执行 THO 需要 design 视图。顶层可能正链接到 Abstract 视图；请按「将 Block 从 Abstract 切换为 Design 视图」一节的说明，把要优化的 block 切换为 design 视图。按下列步骤设置 THO：
\begin{enumerate}
  \item 若在块级 \cmd{compile\_fusion} 之后执行 THO，必须用下列命令把 CCD useful skew balancing offset 以及集成时钟门控单元的 latency offset 传播到顶层时序上下文：
\begin{lstlisting}
fc_shell> promote_clock_data -auto_clock all -latency_offset
fc_shell> promote_clock_data -auto_clock all -balance_points
\end{lstlisting}
  \item 要启用并对 block 执行 THO，确保要优化的 block 已允许编辑：
\begin{lstlisting}
fc_shell> set_editability -blocks {list of blocks} -value true
\end{lstlisting}
  \item 用下列命令验证 block 可编辑性：
\begin{lstlisting}
fc_shell> report_editability -blocks [get_blocks -hierarchical]
\end{lstlisting}
  \item 在 THO 期间，子块内的拥塞更新、route-driven estimation（RDE）计算以及 ECO 布线等操作可用分布式处理运行。用下列命令进行设置。必须按示例指定自己的队列参数（如 \texttt{qsub} 参数）：
\begin{noteBox}
若未指定 host options，工具会在主进程中串行运行各子块操作。
\end{noteBox}
\begin{lstlisting}
fc_shell> set_host_options -name for_tho -num_processes 4 -max_cores 16
 -submit_command [list qsub -P queue -pe mt 16 -l mem_free=100G -cwd]
fc_shell> set_hierarchy_options -host_option for_tho -blocks {list of
 blocks}
\end{lstlisting}
  \item 用 \cmd{init\_hier\_optimization} 启用 THO。按阶段使用不同选项，见后续各节。
\end{enumerate}

\subsection{在 compile\_fusion final\_place 阶段执行 THO}
\subsubsection*{Performing THO During compile\_fusion final\_place Stage}

在 \cmd{compile\_fusion} 的 \texttt{final\_place} 阶段，可跨顶层与各 block 透明地执行布局。布局发生在各自 block 边界内，同时又是全局的——忽略 pin 位置以获得直通路径。布局可由线长、时序与拥塞驱动。

THO 提供基于全局布局细化 pin 位置的选项。可在 THO 的粗布局步骤中，为指定 block 集细化 pin 位置。
\begin{enumerate}
  \item 要细化 pin 位置并允许 pin 移动：
\begin{lstlisting}
fc_shell> set_hierarchy_options -enable_pin_movement true -blocks \
 {list of blocks}
\end{lstlisting}
  \item 要用 \cmd{report\_hierarchy\_options} 报告 \cmd{set\_hierarchy\_options} 所施加的设置。
  \item 要用下列命令以 THO 运行 \texttt{final\_place} 阶段：
\begin{lstlisting}
fc_shell> init_hier_optimization -flow pre_route
fc_shell> compile_fusion -from final_place -to final_place
fc_shell> commit_hier_optimization
\end{lstlisting}
\end{enumerate}

\cmd{init\_hier\_optimization} 会发出类似下列消息，确认哪些 block 可用于优化：
\begin{lstlisting}
Information: Sub-block u1/a/i1 (Reference: ablock.design) is treated
 as modifiable. (HOPT-001)
\end{lstlisting}

全局布局之后，THO 调用 \cmd{place\_pins} 创建新的 pin 布局。

\subsection{在 compile\_fusion final\_opto 阶段执行 THO}
\subsubsection*{Performing THO During compile\_fusion final\_opto Stage}

Fusion Compiler 在 \texttt{final\_opto} 阶段执行透明层次优化，确保所有块间路径与块-顶层路径都像扁平设计一样被优化。这样，CCD 优化可完整处理这些路径上的违规。
\begin{lstlisting}
fc_shell> init_hier_optimization -flow pre_route
fc_shell> compile_fusion -from final_opto -to final_opto
fc_shell> commit_hier_optimization
\end{lstlisting}

\subsection{compile\_fusion 之后重新预算}
\subsubsection*{Rebudgeting After compile\_fusion}

在带 THO 的 \cmd{compile\_fusion} 之后，块接口时序相对各 block 单独实现时会发生变化。由于 block 现已基于顶层时序图景优化，在继续对 block 运行 \cmd{clock\_opt} 之前，必须为 block 创建更新后的时序约束。THO 期间可能已修改 CCD useful skew balance point offset；这些 offset 也需要变成块上下文专用，以便时钟树综合步骤能将其纳入考虑。

下例展示如何执行预算：
\begin{lstlisting}
fc_shell> set budget_instances "u1/a/i1 u1/b/i2 ..."
fc_shell> set_budget_options -reset -all
fc_shell> set_budget_options -add_blocks $budget_instances
fc_shell> set_app_options -list { plan.budget.use_ccd_latency true }
fc_shell> set_app_options -list { plan.budget.write_hold_budgets false }
fc_shell> set_app_options -list { plan.budget.estimate_timing_mode
 false }
fc_shell> compute_budget_constraints -setup_delay -latency_targets actual
 -balance false -boundary
fc_shell> write_budgets -output budgets-post-compile -force -nosplit
 -include_balance_points -include_pin_latency
\end{lstlisting}

\subsection{在 route\_opt 或 hyper\_route\_opt 期间执行 THO}
\subsubsection*{Performing THO During route\_opt or hyper\_route\_opt}

在实现流程末尾执行 THO，可获得最佳优化潜力，因为优化的是完整设计。\cmd{route\_opt} 与 \cmd{hyper\_route\_opt} 对顶层与各 block 一起执行后布线优化。优化后的 ECO 布线会在所有已打开编辑的子块上执行。

THO 优化已完全集成到 \cmd{route\_opt} 与 \cmd{hyper\_route\_opt}。在单主机上使用 \cmd{hyper\_route\_opt} 或 \cmd{route\_opt} 的示例如下：
\begin{lstlisting}
fc_shell> init_hier_optimization -flow post_route
fc_shell> route_opt  or hyper_route_opt
fc_shell> commit_hier_optimization
\end{lstlisting}

使用 \cmd{route\_opt} 时，可在每次 \cmd{route\_opt} 之后并行执行必须进行的 ECO 布线操作，以缩短整体运行时间。须按「设置透明层次优化」一节设置 host options。\cmd{hier\_route\_eco} 在各 block 中执行 ECO 布线，示例如下：
\begin{lstlisting}
fc_shell> init_hier_optimization -flow post_route \
 -detached_hier_route_eco
fc_shell> route_opt
fc_shell> hier_route_eco
fc_shell> commit_hier_optimization
\end{lstlisting}

\begin{noteBox}
以这种方式使用 \cmd{hyper\_route\_opt} 时，工具会忽略该设置，不会分发 ECO 布线。
\end{noteBox}
